
\subsubsection{Interpretation of Null Result}

The hypothesis that CV predicts cross-benchmark stability is refuted. Three explanations emerge:

\textbf{1. Cross-benchmark $\rho$ conflates reliability and validity.} The dependent variable---mean cross-benchmark Spearman $\rho$---was intended to measure ranking stability (reliability), but it also reflects construct divergence (validity). Low $\rho$ could indicate unreliable measurement (noisy benchmarks produce inconsistent rankings) or valid construct divergence (benchmarks measure orthogonal dimensions and should have low $\rho$ by design). CV may predict reliability but cannot predict validity. Since cross-benchmark $\rho$ conflates both, CV shows weak correlation.

\textbf{2. Score variance has multiple sources CV cannot disambiguate.} High CV can arise from measurement noise (inconsistent scoring), valid task heterogeneity (items spanning easy-to-hard difficulty by design), ceiling/floor effects (score clustering at extremes), or construct-specific variance (different trust dimensions have intrinsically different score distributions). CV collapses these sources into a single number and cannot reliably predict which benchmarks are noisy versus well-designed with intentional heterogeneity.

\textbf{3. Simple meta-features are insufficient---multi-feature models needed.} Benchmark quality is multi-faceted (construct validity, reliability, discriminative power, item difficulty calibration). A single statistic (CV) cannot capture this complexity. Future work should test whether features in combination (CV, IQR, split-half reliability, score ceiling proximity, skewness) predict cross-benchmark stability better than CV alone.

\subsubsection{Theoretical Contribution: Construct Divergence}

The cross-benchmark correlation analysis reveals a surprising pattern: negative and near-zero correlations between benchmarks ostensibly measuring ''trustworthiness.'' The strongest negative correlation---FaithfulQA vs. FinTrust ($\rho$=-0.568)---is particularly striking. Both benchmarks claim to measure trust, yet they rank models inversely.

Two interpretations are possible: (1) one or both benchmarks are poorly designed, producing inverted rankings due to noise or bias (methodological failure), or (2) FaithfulQA measures hallucination detection (factual grounding) while FinTrust measures financial ethical reasoning (value alignment), which are orthogonal dimensions within the broader ''trust'' construct (construct divergence). Psychometric theory supports interpretation 2. Campbell \& Fiske (1959) distinguish convergent validity (measures of the same construct should correlate positively) from discriminant validity (measures of different constructs should correlate weakly or negatively). Negative $\rho$ suggests good discriminant validity---they measure distinct sub-dimensions.

This finding challenges a common assumption in benchmark meta-analysis: that benchmarks within a domain (e.g., ''trust evaluation'') measure overlapping constructs and should exhibit positive cross-benchmark correlations. Our results suggest an alternative: disagreement may reflect valid multi-dimensionality, not measurement failure. Before interpreting low cross-benchmark $\rho$ as a quality issue, researchers should validate construct structure via factor analysis and test external validity against criteria (correlation with human trust judgments, real-world deployment failures).

\subsubsection{Limitations}

\textbf{Mock data validity threat (CRITICAL).} This analysis uses a mock benchmark corpus with synthetic model scores, not real scraped leaderboards from TrustLLM, HaluBench, or TruthfulQA. Real leaderboards have structured biases (model selection bias, evaluation protocol heterogeneity, temporal effects) absent from mock data. Mock data generation may have unintentionally created null relationships (e.g., random CV-$\rho$ pairing), making the r=-0.486 finding an artifact rather than a true empirical pattern. The negative cross-benchmark correlations (FaithfulQA-FinTrust $\rho$=-0.568) may also be mock data artifacts. Rerunning h-e1 with real leaderboard data is required before accepting the null result as valid. Current conclusion validity is uncertain---null result may be valid, or it may be a data artifact.

\textbf{Domain restriction.} The hypothesis was tested exclusively on trust evaluation benchmarks. The null result applies only to this domain. Math/reasoning benchmarks (GSM8K, MATH) may have high CV by design (difficulty stratification from easy arithmetic to olympiad problems), where high CV indicates good task coverage, not instability. Vision-language benchmarks (VQA, COCO) and general NLP (GLUE, SuperGLUE) may exhibit different CV-stability dynamics. Do not generalize this null result beyond trust evaluation.

\textbf{CV operationalization for multi-dimensional benchmarks.} TrustLLM reports 8 sub-dimension scores (truthfulness, safety, fairness, robustness, privacy, ethics, machine ethics, toxicity), not a single overall score. CV computation method (averaging across dimensions versus using a single dimension) was not documented in validation reports. If one dimension has high CV but others have low CV, averaging masks dimension-specific patterns. Alternative operationalizations (average CV, max CV, dimension-specific) were not tested for sensitivity.

\textbf{Mechanism validation blocked.} The gate-driven workflow blocked mechanism hypotheses when the existence hypothesis failed MUST\_WORK criteria. We cannot distinguish whether the hypothesis is fundamentally false (CV genuinely doesn't predict stability via any mechanism) or whether the mechanism is wrong (CV predicts stability, but not via the hypothesized pathway). Without mechanism testing, we don't know why CV fails.

\subsubsection{Broader Impact}

Positive impacts include providing methodological rigor (pre-registration, power analysis, falsification framework) as a template for hypothesis-driven benchmark research, contributing negative evidence that constrains future tool development toward multi-feature or construct-aware approaches, and raising construct validity awareness. Risks include over-simplification (simple CV thresholds may mislead researchers if adopted without understanding limitations) and premature dismissal (this null result for CV doesn't mean all leaderboard meta-features fail).
